\documentclass{article}
\usepackage{amsmath, amssymb, amsthm}
\usepackage{hyperref}
\usepackage{geometry}
\geometry{margin=1in}

\title{Erd\H{o}s Problem \#995}
\date{Last updated: 2025-09-07}
\author{Source: erdosproblems.com}

\begin{document}
\maketitle

\noindent\textbf{Status:} OPEN \quad \textbf{No prize}

\noindent\textbf{Tags:} analysis, discrepancy

\medskip

\noindent\textbf{Problem Statement:}

Let $n_1&#60;n_2&#60;\cdots$ be a lacunary sequence of integers and $f\in L^2([0,1])$. Estimate the growth of, for almost all $\alpha$,\[\sum_{1\leq k\leq N}f(\{ \alpha n_k\}).\]For example, is it true that, for almost all $\alpha$,\[\sum_{1\leq k\leq N}f(\{ \alpha n_k\})=o(N\sqrt{\log\log N})?\]




Erd\H{o}s \cite{Er49d} constructed a lacunary sequence and $f\in L^2([0,1])$ such that, for every $\epsilon&#62;0$, for almost all $\alpha$\[\limsup_{N\to \infty}\frac{1}{N(\log\log N)^{\frac{1}{2}-\epsilon}}\sum_{1\leq k\leq N}f(\{\alpha n_k\})=\infty.\]Erd\H{o}s also proved that, for every lacunary sequence and $f\in L^2$, for every $\epsilon&#62;0$, for almost all $\alpha$,\[\sum_{1\leq k\leq N}\sum_{1\leq k\leq N}f(\{\alpha n_k\})=o( N(\log N)^{\frac{1}{2}+\epsilon}).\]Erd\H{o}s \cite{Er64b} thought that his lower bound was closer to the truth.




References

[Er49d] Erd\H{o}s, P., On the strong law of large numbers . Trans. Amer. Math. Soc. (1949), 51--56.

[Er64b] Erd\H{o}s, P., Problems and results on diophantine approximations . Compositio Math. (1964), 52-65.


\bigskip
\noindent\small{Source: \url{https://www.erdosproblems.com/995}}
\end{document}
